\expandafter\let\csname Bs\endcsname\undefined
\expandafter\let\csname CP\endcsname\undefined
\expandafter\let\csname CX\endcsname\undefined
\expandafter\let\csname CZ\endcsname\undefined
\expandafter\let\csname Es\endcsname\undefined
\expandafter\let\csname Ps\endcsname\undefined
\expandafter\let\csname QFT\endcsname\undefined
\expandafter\let\csname Ss\endcsname\undefined
\expandafter\let\csname Vs\endcsname\undefined
\expandafter\let\csname Ws\endcsname\undefined
\expandafter\let\csname bits\endcsname\undefined
\expandafter\let\csname class\endcsname\undefined
\expandafter\let\csname poly\endcsname\undefined
\expandafter\let\csname qbig\endcsname\undefined
\expandafter\let\csname qcircuit\endcsname\undefined
\expandafter\let\csname qclink\endcsname\undefined
\expandafter\let\csname qcross\endcsname\undefined
\expandafter\let\csname qcwire\endcsname\undefined
\expandafter\let\csname qdot\endcsname\undefined
\expandafter\let\csname qg\endcsname\undefined
\expandafter\let\csname qgb\endcsname\undefined
\expandafter\let\csname qlab\endcsname\undefined
\expandafter\let\csname qlink\endcsname\undefined
\expandafter\let\csname qmeter\endcsname\undefined
\expandafter\let\csname qopen\endcsname\undefined
\expandafter\let\csname qrlab\endcsname\undefined
\expandafter\let\csname qtarg\endcsname\undefined
\expandafter\let\csname qtext\endcsname\undefined
\expandafter\let\csname qwire\endcsname\undefined
\newcommand{\class}[1]{\mathsf{#1}}
\newcommand{\QFT}{\operatorname{QFT}}
\newcommand{\poly}{\operatorname{poly}}
\newcommand{\Vs}{\mathcal{V}}
\newcommand{\Ws}{\mathcal{W}}
\newcommand{\Es}{\mathcal{E}}
\newcommand{\Ss}{\mathcal{S}}
\newcommand{\Bs}{\mathcal{B}}
\newcommand{\Ps}{\mathcal{P}}
\newcommand{\CX}{\mathrm{CX}}
\newcommand{\CZ}{\mathrm{CZ}}
\newcommand{\bits}[1]{\mathbb{B}_{#1}}
\newcommand{\CP}{\mathbb{CP}}





\tikzset{
  qw/.style={line width=0.5pt},
  qgt/.style={draw, line width=0.5pt, fill=white, rounded corners=0.5pt,
              minimum width=6.5mm, minimum height=6.5mm, inner sep=1.5pt, font=\small},
  qcd/.style={circle, fill=black, minimum size=1.6mm, inner sep=0pt}
}
\newcommand{\qcircuit}[1]{\begin{tikzpicture}[x=1.05cm,y=0.85cm,baseline]#1\end{tikzpicture}}
\newcommand{\qwire}[3]{\draw[qw] (#2,#1) -- (#3,#1);}
\newcommand{\qcwire}[3]{\draw[qw,double,double distance=1.1pt] (#2,#1) -- (#3,#1);}
\newcommand{\qlink}[3]{\draw[qw] (#1,#2) -- (#1,#3);}
\newcommand{\qclink}[3]{\draw[qw,double,double distance=1.1pt] (#1,#2) -- (#1,#3);}
\newcommand{\qg}[3]{\node[qgt] at (#1,#2) {$#3$};}
\newcommand{\qgb}[4]{\node[qgt,minimum width=#4] at (#1,#2) {$#3$};}
\newcommand{\qdot}[2]{\node[qcd] at (#1,#2) {};}
\newcommand{\qopen}[2]{\draw[qw,fill=white] (#1,#2) circle[radius=1.1mm];}
\newcommand{\qtarg}[2]{%
  \draw[qw,fill=white] (#1,#2) circle[radius=2.1mm];
  \draw[qw] (#1,#2) ++(-2.1mm,0) -- ++(4.2mm,0);
  \draw[qw] (#1,#2) ++(0,-2.1mm) -- ++(0,4.2mm);}
\newcommand{\qcross}[2]{%
  \draw[qw] (#1,#2) ++(-1.5mm,-1.5mm) -- ++(3mm,3mm);
  \draw[qw] (#1,#2) ++(-1.5mm,1.5mm) -- ++(3mm,-3mm);}
\newcommand{\qmeter}[2]{%
  \node[qgt,minimum width=7.5mm] at (#1,#2) {};
  \draw[qw] (#1,#2) ++(-2.2mm,-1.3mm) arc[start angle=180,end angle=0,radius=2.2mm];
  \draw[qw] (#1,#2) ++(0,-1.3mm) -- ++(1.7mm,2.9mm);}
\newcommand{\qlab}[3]{\node[left,inner sep=3pt,font=\small] at (#1,#2) {$#3$};}
\newcommand{\qrlab}[3]{\node[right,inner sep=3pt,font=\small] at (#1,#2) {$#3$};}
\newcommand{\qtext}[3]{\node[font=\small] at (#1,#2) {#3};}
\newcommand{\qbig}[5]{%
  \filldraw[qw,fill=white,rounded corners=0.5pt] (#1-#4,#2+0.32) rectangle (#1+#4,#3-0.32);
  \node[font=\small] at (#1,{(#2+#3)/2}) {$#5$};}
